В прошлом билете мы выяснили, что $(x + y, x + \omega y) = (x + y, x + \omega^2 y) = (x + \omega^2y, x + \omega y) = \lambda$. Если
 вынести $\lambda$, то мы получим взаимно простые числа, дающие в произведении $z^3$. Следовательно, в разложении на простые каждый сомножитель присутствует в степени, делящейся на 3. Поэтому данное произведение можно представить в следующем виде: $z^3 = (x + y)(x + \omega y)(x + \omega^2y) = ([\lambda]u_1a^3)([\lambda]u_2b^3)([\lambda]u_3c^3)$, где $u_1, u_2, u_3$ - обратимые элементы, $[\lambda]$ = $\lambda$ или 1.
Если бы выполнялось равенство $a^3 + b^3 + c^3 = 0$, то мы получили бы решение, которое «меньше» исходного (к примеру, ясно,
что норма a, b, c меньше нормы z). Тогда бы мы могли прийти к
противоречию рассмотрев тройку решения с минимальным наи
большим значением N(x), N(y), N(z). Такой метод называется методом спуска. Но сразу это равенство получить не удастся.
Действительно, сумма $(x + y) + (x + \omega y) + (x + \omega^2y) \neq 0$. 
\\
Однако, используя равенство $1 + \omega + \omega^2 = 0$ легко получить, что $(x+y)+\omega(x+\omega y)+\omega^2(x+\omega^2y) = (x+\omega x+\omega^2x)+(y+\omega^2y+\omega^4y) = 0$.\\
Следовательно,
\begin{center}
    $[\lambda](u_1a^3 + \omega u_2b^3 + \omega^2 u_3c^3)$ = 0.
\end{center}
Так как $\omega, \omega^2$ - обратимые элементы, а на $u_1 и \lambda$ можно сократить, мы получаем равенство
\begin{center}
    $a^3 +  u_1'b^3 +
     \omega^2 u_2'c^3 = 0$
\end{center}
Это "почти" то, что нам нужно. От одного лишнего коэффициента можно избавиться, рассмотрев более общее уравнение
$x^3 + y^3 = u z^3$, где u - обратимый элемент.
\\
Для того, чтобы показать, что $u_1' = \pm1$, надо проверить, что сумма  $a^3 +  u_1'b^3$ будет кубом (с точностью до обратимого элемента)
только в том случае, если $u_1' = \pm1$
\\
\\
\Lemma
	Если $r_xx^3 + r_xy^3 = r_zz^3\lambda^{3k}$ для некоторых $x, y, z \in \Z[\omega]$, $r_x, r_y, r_z \in \Z[\omega]^*$, $k \in \N$, причем $\lambda \nmid xyz$, то выполнено следующее:
	\begin{enumerate}
		\item $k \ge 2$
		\item $r_x = \pm r_y$
	\end{enumerate}


\Proof
	Рассмотрим равенство из условия по модулю 9: $\pm r_x \pm r_y \equiv_9 \pm r_z\lambda^{3k}$. Если $k = 1$, то $0 \le N(\pm r_x \pm r_y) \le 2 < N(r_z\lambda^{3})< 9$, поэтому данное сравнение не может быть выполнено. Если же $k \ge 2$, то правая часть делится на 9, поэтому $\pm r_x \pm r_y = 0$, откуда $r_x = \pm r_y$.
\EndProof
